\chapter{Altair Guide}
A good guide on how to use Altair to visualize Pandas \lstinline|DataFrame| objects, can be found on \url{https://altair-viz.github.io/user_guide/data.html}.

Another great guide is found on \url{https://www.geeksforgeeks.org/python/introduction-to-altair-in-python/}.

\section{"My" Guide on Altair}
In Altair, we can plot Pandas \lstinline|DataFrame| objects with \lstinline|$plt$ = altair.Chart($D$).mark_$m$().encode($ch_1$=$co_1$, ...)|---where \(D\) is a \lstinline|DataFrame|, \(m\) is a mark, and \lstinline|$ch_1$=$co_1$| specifies which channels are expressing which attribute or column of the table.

Marks are specified with the methods listed in \autoref{tab:1}.
Possible channels are the parameters of the \lstinline|encode| method listed in \autoref{tab:2}.
We can also hint which data type \(T\) a column \(co\) has by passing \lstinline|$co$:$T$| to some channel.
\(T\) can be a value of \autoref{tab:3}.
Type hinting may be necessary because Altair needs to know how to interpret the data and determine e.g., scales from that.
To combine and layer two generated charts, we can just write \lstinline|$plt_1$ + $plt_2$|.

\begin{table}[t]
    \centering
    \begin{tabular}{lll}
        \hline
        \textbf{Mark} & \textbf{Method} & \textbf{Description} \\
        \hline
        Point   & \texttt{.mark\_point()}   & Scatter plot points \\
        Circle  & \texttt{.mark\_circle()}  & Like point but size-scaled \\
        Line    & \texttt{.mark\_line()}    & Line chart \\
        Bar     & \texttt{.mark\_bar()}     & Bar chart \\
        Area    & \texttt{.mark\_area()}    & Area chart \\
        Tick    & \texttt{.mark\_tick()}    & Small tick marks \\
        Rect    & \texttt{.mark\_rect()}    & Heatmaps \\
        Rule    & \texttt{.mark\_rule()}    & Reference lines \\
        Text    & \texttt{.mark\_text()}    & Labels \\
        Square  & \texttt{.mark\_square()}  & Square markers \\
        \hline
    \end{tabular}
    \caption{Marks their methods in Altair.}
    \label{tab:1}
\end{table}

\begin{table}[t]
    \centering
    \begin{tabular}{ll}
        \hline
        \textbf{Channel} & \textbf{Description} \\
        \hline
        \texttt{x}        & Horizontal position \\
        \texttt{y}        & Vertical position \\
        \texttt{color}    & Color of mark \\
        \texttt{size}     & Size of mark \\
        \texttt{shape}    & Shape of mark \\
        \texttt{opacity}  & Transparency \\
        \texttt{tooltip}  & Hover information \\
        \texttt{row}      & Facet by rows \\
        \texttt{column}   & Facet by columns \\
        \texttt{detail}   & Additional grouping \\
        \texttt{order}    & Drawing order \\
        \texttt{text}     & Text labels \\
        \hline
    \end{tabular}
    \caption{Channels in Altair.}
    \label{tab:2}
\end{table}

\begin{table}[t]
    \centering
    \begin{tabular}{ll}
        \hline
        \textbf{Type} & \textbf{Description} \\
        \hline
        \texttt{Q} & Quantitative \\
        \texttt{N} & Nominal \\
        \texttt{O} & Ordinal \\
        \texttt{T} & Temporal \\
        \hline
    \end{tabular}
    \caption{Possible typehints for columns in Altair.}
    \label{tab:3}
\end{table}
